\documentclass[11pt]{article}
\usepackage[utf8]{inputenc}
\usepackage[T1]{fontenc}
\usepackage{amsmath,amssymb,amsthm}
\usepackage{booktabs}
\usepackage{graphicx}
\usepackage{xcolor}
\usepackage{microtype}
\usepackage[margin=1in]{geometry}
\usepackage[hidelinks]{hyperref}
\usepackage{natbib}
\usepackage{enumitem}

\newtheorem{lesson}{Lesson}

\title{\bf Can Bayes Help Information Retrieval for LLM Systems?\\
\large Structure and Adaptation, Not Uncertainty: A Controlled Audit}
\author{Anonymous}
\date{\today}

\begin{document}
\maketitle

\begin{abstract}
Bayesian and ``uncertainty'' language pervades retrieval-augmented generation (RAG) and
text-to-SQL, usually without a controlled check that it beats the trivial baselines practitioners
already use. We run that check. Treating retrieval as a clean, executable testbed---schema linking
for text-to-SQL, plus multi-hop and single-hop document retrieval---we instrument the five places a
Bayesian idea can enter a retriever (calibrated relevance/uncertainty, signal fusion, structured
joint selection, set-level diversity, and online adaptation) and compare each against simple
heuristics (cosine top-$k$, reciprocal-rank fusion, a max-excluded-similarity confidence, fixed
top-$k$ context). Across six datasets (BIRD, Spider, BEAVER, HotpotQA, 2WikiMultiHopQA, SciFact) the
verdict is consistent and, we argue, generalizable. Bayes helps through exactly two mechanisms:
\emph{structure}---a graph prior over the jointly-relevant subset, which improves multi-hop recall
and downstream execution accuracy---and \emph{adaptation}---online posterior updating, which helps in
correlated, repeated workloads. It helps \emph{only} when those conditions hold: a structural prior is
neutral-to-harmful on single-hop queries and \emph{catastrophic} when imposed without genuine
connectivity, and adaptation requires a workload that repeats. Crucially, Bayes \emph{fails} where it
is most often invoked---as an uncertainty or decision signal---losing to one-line baselines at both
abstention and context-budget selection. And where structure does help, it is the \emph{graph}, not
the \emph{Bayes}: a shortest-path heuristic and a PageRank diffusion match the full subgraph
posterior everywhere we test. We distill four lessons and a practitioner's decision guide, and
identify the one place the full Bayesian object earns its keep: marginalization auto-gates structure
across heterogeneous workloads.
\end{abstract}

\section{Introduction}

Retrieval is now a load-bearing component of large-language-model (LLM) systems. RAG pipelines
retrieve passages before generating; text-to-SQL systems retrieve the relevant tables and columns
(``schema linking'') before writing a query. In both, the literature is saturated with Bayesian and
uncertainty-quantification (UQ) framing: posterior relevance, calibrated confidence, principled
fusion of signals, diversity priors, abstention under uncertainty. Yet the field rarely asks the
control question a statistician would ask first: \emph{does any of this beat the trivial baseline?}
Cosine top-$k$, reciprocal-rank fusion, and a single held-out threshold on a transparent score are
strong, cheap, and ubiquitous. If a Bayesian apparatus does not beat them, it is decoration.

This article is that controlled audit. We are deliberately willing to report a boring answer. Our
headline is nuanced rather than triumphant---and we think that is exactly the contribution. The
answer is: \textbf{Bayes helps retrieval through \emph{structure} and \emph{adaptation}, but not
through \emph{uncertainty} or \emph{decision} signals,} and each positive mechanism helps only under
identifiable conditions that we map precisely.

\paragraph{Why text-to-SQL, and why a multi-dataset sweep.} Text-to-SQL retrieval is an unusually
clean testbed: the corpus (database tables) is finite and known, the gold relevant set is recoverable
from the gold SQL, there is a \emph{true} relational graph (foreign keys), and success is checkable by
execution. That lets us measure not just retrieval recall but downstream execution accuracy. We then
test whether each finding generalizes to open-domain RAG by replicating on multi-hop QA (HotpotQA,
2WikiMultiHopQA), single-hop retrieval (SciFact), and a correlated enterprise warehouse (BEAVER). The
result is a $5\times5$ matrix---five Bayesian entry points by five retrieval regimes---that pins down
\emph{when} each mechanism helps.

\paragraph{Contributions.}
\begin{itemize}[leftmargin=1.4em,itemsep=2pt]
\item A taxonomy of the five places a Bayesian idea enters an LLM retriever, each paired with the
trivial baseline it must beat (\S\ref{sec:taxonomy}).
\item A controlled, reproducible audit across six datasets and both domains (SQL, RAG), with the
master evidence table as its spine (\S\ref{sec:results}).
\item The central empirical finding: a \emph{connectivity boundary}. Structure's value is monotone in
how connected the evidence is---negative on single-hop, growing with hop count and corpus
correlation---and an imposed graph without real connectivity is actively harmful
(\S\ref{sec:boundary}).
\item A constructive positive: an adaptive, topology-routed retriever that beats every fixed prior,
and the observation that the subgraph posterior's marginalization auto-gates structure across
heterogeneous workloads (\S\ref{sec:constructive}).
\item Four lessons and a practitioner's decision guide (\S\ref{sec:lessons}--\S\ref{sec:guidance}).
\end{itemize}

\section{Background: ``Bayesian retrieval'' already exists}\label{sec:background}

It is worth stating plainly that probabilistic and Bayesian retrieval is not new; much of what is
proposed as ``Bayesian retrieval for LLMs'' reinvents a mature lineage. The probabilistic relevance
framework \citep{robertson1976relevance} and its modern form BM25 \citep{robertson2009bm25} already
\emph{is} a probability of relevance. Language-model IR \citep{ponte1998lm} with Dirichlet smoothing
\citep{zhai2004smoothing} is a Bayesian predictive distribution over terms. Risk-minimization and
decision-theoretic IR \citep{zhai2006risk} cast retrieval as expected-loss minimization, and diversity
via determinantal point processes \citep{kulesza2012dpp} is an explicitly probabilistic repulsion
model. Our question is therefore narrow and modern: granting that classical probabilistic IR works,
does \emph{LLM-era} Bayes---posteriors over LLM outputs, graph priors over retrieved sets,
posterior-as-uncertainty---add anything beyond cheap heuristics? That is what we instrument.

\section{A taxonomy: five places Bayes can enter an LLM retriever}\label{sec:taxonomy}

We organize the audit by where a Bayesian object could plausibly do work, and pair each with the
baseline it must beat.

\begin{enumerate}[leftmargin=1.6em,itemsep=3pt]
\item \textbf{Calibrated relevance / uncertainty.} A posterior probability of relevance or of answer
correctness, used to rank or to abstain. \emph{Baseline:} a held-out threshold on a transparent
signal (max excluded similarity; a verifier).
\item \textbf{Signal fusion.} A Bayesian combination of heterogeneous scores (lexical, dense, value
overlap). \emph{Baseline:} reciprocal-rank fusion (RRF) or plain cosine.
\item \textbf{Structured joint selection.} A prior over the jointly-relevant \emph{subset}, coupling
items through a relational graph (foreign keys, hyperlinks, entity links). \emph{Baseline:} independent
top-$k$; a shortest-path connectivity heuristic.
\item \textbf{Set-level diversity.} A repulsion prior (DPP) to avoid redundant retrievals.
\emph{Baseline:} maximal marginal relevance (MMR).
\item \textbf{Online adaptation.} A prior updated from feedback as queries arrive. \emph{Baseline:} a
static retriever.
\end{enumerate}

For each, the natural Bayesian object is well defined. The empirical question is whether it beats the
baseline, and under what conditions.

\section{Experimental setup}\label{sec:setup}

\paragraph{Datasets and regimes.} We span five retrieval regimes:
(S1) \emph{orthogonal SQL}---BIRD and Spider, whose databases cover unrelated domains;
(S2) \emph{correlated SQL}---the BEAVER enterprise warehouse \citep{chen2024beaver} (97 tables sharing
prefixes and keys, mean 4 gold tables per query);
(S3) \emph{single-hop}---single-table SQL queries and SciFact \citep{wadden2020scifact} claim
verification (5{,}183 abstracts, $\approx$1 gold doc);
(S4) \emph{multi-hop RAG}---HotpotQA distractor \citep{yang2018hotpotqa} (10 passages, 2 gold);
(S5) \emph{graph RAG}---2WikiMultiHopQA \citep{ho2020twowiki} with four reasoning types.

\paragraph{Methods.} For structured selection we build, per query, a cross-fit logistic
\emph{unary} relevance score over interpretable features (dense cosine, BM25, name/column overlap,
max column cosine, value-token overlap), then couple items through the relational graph in three ways:
(i) an Ising/MRF posterior over connected subsets with exact marginalization where tractable
($\sum_{\text{subsets}} \exp(\sum_i a_i x_i + \beta\sum_{(i,j)\in E} x_i x_j)$), (ii) personalized
PageRank diffusion seeded by the unary scores, and (iii) a shortest-path foreign-key closure
heuristic. Hyperparameters ($\beta$, diffusion $\alpha$, closure bonus $\gamma$) are chosen on held-out
queries. Diversity uses MMR and a $k=2$ DPP; UQ uses the LLM's softmax/posterior versus a
max-excluded-similarity signal; adaptation uses online Beta-Bernoulli popularity and a naive-Bayes
term$\rightarrow$table model updated prequentially. Baselines are cosine top-$k$, BM25, RRF, and
fixed-$k$ context. We report recall at $|{\rm gold}|$, downstream execution accuracy (EX) for SQL, and
AUROC for abstention; confidence intervals are paired bootstrap. All embeddings are
\texttt{text-embedding-3-small}; generation/verification use small chat models. Total API spend was
$\approx\$6$; all scripts and cached data are released.

\section{Results, by taxonomy cell}\label{sec:results}

Table~\ref{tab:master} is the spine of the paper. We summarize the cells in prose, then devote
\S\ref{sec:boundary}--\S\ref{sec:constructive} to the two findings that organize everything.

\begin{table}[t]
\centering\small
\caption{Master evidence table. Each row is a (Bayesian object, baseline) contest in a given regime.
Bold verdicts are the load-bearing results. Numbers are recall at $|{\rm gold}|$ unless noted.}
\label{tab:master}
\begin{tabular}{@{}p{3.5cm}p{3.3cm}cp{4.5cm}@{}}
\toprule
\textbf{Cell / regime} & \textbf{Bayesian object vs.\ baseline} & \textbf{Verdict} & \textbf{Key numbers}\\
\midrule
Structured selection, SQL (S1) & FK-graph MRF vs.\ cosine; vs.\ FK-closure & \textbf{Bayes/structure helps} & MRF 0.805/0.822/0.787 vs.\ cosine 0.720/0.673/0.678; vs.\ closure $+$0.02--0.04; \textbf{EX $+$5.7pp} [2.1,9.4]\\
Structured selection, correlated SQL (S2) & join-graph diffusion/MRF vs.\ cosine & \textbf{helps more} & cosine craters 0.425; PageRank 0.550, MRF 0.556 ($+$12--14pp)\\
Structured selection, multi-hop RAG (S4) & passage-link graph vs.\ cosine & \textbf{helps (generalizes)} & MRF$-$cosine $+$0.104 [.090,.118]; bridge $+$17pp\\
Graph RAG by reasoning type (S5) & entity-graph PageRank/MRF vs.\ cosine & \textbf{helps chained, hurts independent} & bridge $+$0.21, compositional $+$0.17, inference $+$0.13, \emph{comparison} $-$0.22; PageRank$\approx$MRF (0.805$\approx$0.804)\\
Single-hop control, SQL (S3) & FK diffusion/MRF vs.\ cosine & \textbf{structure hurts / MRF neutral} & PageRank 0.727 $<$ cosine 0.839 ($-$0.11); MRF 0.857\\
Single-hop control, RAG (S3) & cosine-kNN graph diffusion vs.\ cosine & \textbf{structure destructive} & diffusion 0.016 vs.\ cosine 0.608 ($-$0.59)\\
\midrule
Online adaptation, correlated (S2) & online naive-Bayes / popularity vs.\ static cosine & \textbf{Bayes wins} & naive-Bayes 0.510 vs.\ cosine 0.425 ($+$8.5pp); learning curve $+$0.07\\
\midrule
Decision / context budget (S1) & posterior-threshold context vs.\ fixed-$k$; oracle & \textbf{opportunity real, rule fails} & oracle 0.562@2.2 tbl $>$ full 0.500@10; MRF-threshold 0.411 under-retrieves $<$ fixed-$k$ 0.495\\
Calibrated relevance / UQ & posterior vs.\ cosine max-out & \textbf{Bayes loses} & abstention AUROC 0.700 vs.\ 0.763\\
Signal fusion (clean text) & learned fusion vs.\ cosine/RRF & \textbf{no win} & fusion 0.734 $<$ cosine 0.778\\
Set-level diversity & DPP/MMR vs.\ cosine & \textbf{conditional} & $+$7pp on comparison, $0$ on bridge; topology-routing $\approx$0.76\\
\bottomrule
\end{tabular}
\end{table}

\paragraph{Structured selection helps---across both domains.} On orthogonal SQL (S1) the FK-graph
posterior lifts recall well above cosine and, more importantly, lifts downstream execution accuracy by
$5.7$ points [2.1, 9.4]: recovering cosine-invisible ``bridge'' tables (join tables whose names share
no tokens with the question) yields correct SQL. The same mechanism generalizes to multi-hop RAG (S4),
where a passage-link graph beats cosine by $+0.104$ overall and $+0.17$ on bridge questions, and to a
second multi-hop dataset (S5).

\paragraph{Adaptation helps---in repeated workloads.} Streaming the BEAVER queries and updating an
online naive-Bayes term$\rightarrow$table model prequentially beats static cosine by $8.5$ points,
with a genuine learning curve (recall rises from 0.455 early to 0.560 mid-stream). This is the most
textbook-Bayesian win in the study---a prior updated to a posterior as data arrive---and it appears
exactly where the corpus repeats.

\paragraph{Uncertainty and decision signals lose.} Using the LLM's posterior to abstain underperforms
a max-excluded-similarity threshold (AUROC 0.700 vs.\ 0.763). Using the subgraph posterior's marginals
to choose a variable-size context \emph{under}-retrieves and drops execution accuracy to 0.411, below a
fixed top-5 (0.495)---even though the \emph{opportunity} is real, since the oracle gold context (2.2
tables) beats the full schema (10 tables) by 6 points. The posterior is simply not calibrated for a
threshold cutoff; the same failure mode appears at abstention and at budgeting, unifying the two.

\paragraph{Fusion and diversity are not free wins.} On clean text a learned fusion does not beat cosine;
diversity helps only on independent-evidence (``comparison'') queries and is neutral or harmful on
connected-evidence (``bridge'') queries---which sets up the constructive result below.

\section{The connectivity boundary}\label{sec:boundary}

The single most organizing finding is that structure's value is a monotone function of how
\emph{connected} the relevant evidence is.

\begin{itemize}[leftmargin=1.4em,itemsep=2pt]
\item \textbf{Single-hop: structure hurts.} On single-table SQL queries, PageRank diffusion
\emph{lowers} recall@1 from 0.839 to 0.727 ($-0.11$): with one relevant table, diffusion demotes the
lone gold toward its graph neighbors. On single-hop SciFact the effect is extreme---an imposed
cosine-kNN similarity graph collapses recall from 0.608 to 0.016 ($-0.59$), because diffusion spreads
mass into dense clusters of topically-similar-but-irrelevant abstracts and buries the answer. Imposing
structure where there is no genuine relevance-coupling is not merely unhelpful; it is destructive.
\item \textbf{Multi-hop: structure helps, and more so the deeper the chain.} On 2WikiMultiHopQA the
gain rises with reasoning depth and connectivity: $+0.13$ (inference), $+0.17$ (compositional),
$+0.21$ (bridge-comparison), while the one \emph{independent}-evidence type (comparison) goes the other
way, $-0.22$. The comparison type is a within-dataset negative control: same corpus, same method,
opposite sign, determined entirely by whether the evidence is connected.
\item \textbf{Correlation amplifies it.} On the correlated BEAVER warehouse, cosine craters to 0.425
(versus $\approx$0.72 on orthogonal BIRD) and structure's advantage grows to $+$12--14pp. The harder
the corpus is for similarity, the more a true relational graph is worth.
\end{itemize}

\paragraph{It is the graph, not the Bayes.} Wherever structure helps, a personalized-PageRank
diffusion and a shortest-path closure heuristic match the full subgraph posterior: PageRank
$\approx$ MRF on BIRD, BEAVER, HotpotQA, and 2WikiMultiHopQA (e.g., 0.805 vs.\ 0.804 on S5). The
load-bearing ingredient is injecting the \emph{correct relational structure}, not the particular
probabilistic machinery. A practitioner should reach for the cheap heuristic first.

\section{Constructive payoff: route the prior to the query}\label{sec:constructive}

The audit is not only negative. Because the right bias is query-dependent---connectivity for
connected-evidence queries, diversity for independent-evidence queries---a cheap classifier that reads
the query can route between them. On HotpotQA a lexical logistic classifier predicts bridge
vs.\ comparison at 0.936 accuracy; routing structure to bridge and diversity to comparison reaches
recall@2 of 0.748, beating the best \emph{fixed} prior (0.703) by $+0.044$ [.036, .053] and nearly
matching the oracle (0.761). The recommendation is not ``use a graph prior'' but \emph{match the prior
to the query's evidence topology.}

\paragraph{Where the full posterior earns its keep.} There is one place the Bayesian object beats the
heuristic. On a \emph{mixed} workload of single- and multi-hop SQL queries, the subgraph posterior's
marginalization auto-gates: it degrades gracefully on single-hop (where hard diffusion hurts) and wins
on multi-hop, with no predictor. Always-MRF scores 0.803 versus always-PageRank 0.751 and cosine
0.744; an explicit hop-count gate cannot recover this because hop count is hard to predict (classifier
0.799, versus 0.936 for RAG topology). Marginalization buys robustness across heterogeneous query
types precisely when those types are hard to classify---the one genuinely load-bearing use of the full
Bayesian apparatus we found.

\section{Four lessons}\label{sec:lessons}

\begin{lesson}[Structure, not uncertainty]
Bayes helps when it injects \emph{correct structure} a point estimate ignores---here, foreign-key
connectivity that recovers cosine-invisible bridge tables. It is the graph, not the Bayes; the
embedding-similarity graph is the \emph{wrong} structure (connectors are dissimilar), so a
correlation prior does not substitute, and on single-hop data it is actively harmful.
\end{lesson}

\begin{lesson}[The LLM distribution is a proposal, not a posterior]
LLM output distributions are mode-collapsed (near-certain even on ambiguous inputs), so calibrated
uncertainty cannot be read off them. Abstention should use a separate, transparent signal; simple
signals win. The same miscalibration sinks posterior-as-decision (context budgeting), unifying the
uncertainty and decision failures.
\end{lesson}

\begin{lesson}[Adaptation is the other real win, and it is conditional]
Online posterior updating beats a static retriever---but only in correlated, repeated workloads where
there is structure across queries to learn. On one-shot orthogonal benchmarks there is nothing to
adapt to.
\end{lesson}

\begin{lesson}[Marginalization buys robustness, not accuracy]
The one place the full posterior beats the heuristic is heterogeneous workloads: marginalizing over
subset configurations auto-gates structure on and off without a predictor. Where query types are easy
to classify, a cheap explicit router matches it.
\end{lesson}

\section{Practitioner guidance}\label{sec:guidance}

\begin{itemize}[leftmargin=1.4em,itemsep=2pt]
\item \textbf{Reach for graph-aware retrieval} when corpora are large, queries are multi-hop, and a
\emph{true} relational graph (foreign keys, citations, hyperlinks, entity links) aligned with relevance
is available. Expect larger gains as the graph densifies and the corpus correlates.
\item \textbf{Do not} impose structure when queries are single-hop or the only ``graph'' is embedding
similarity---it ranges from useless to destructive.
\item \textbf{Start with the cheap heuristic} (shortest-path closure or PageRank diffusion); only reach
for the subgraph posterior on heterogeneous workloads where graceful auto-gating matters.
\item \textbf{Do not trust posterior uncertainty for abstention or budgeting}; calibrate a held-out
threshold on a transparent signal instead.
\item \textbf{Adapt online} when the workload repeats; route the prior to the query's evidence topology
when it does not.
\end{itemize}

\section{Scope and limitations}\label{sec:limits}

Our own experiments are retrieval-centric and lean on text-to-SQL, where the gold set and the graph are
unambiguous; we treat open-domain RAG through HotpotQA, 2WikiMultiHopQA, and SciFact rather than
web-scale corpora, and we mark SQL-specific claims as such. The adaptation result uses simulated
feedback (gold revealed after each query) on a modest stream (120 queries); it demonstrates the
mechanism and learning curve but not large-scale online deployment. UQ is evaluated with black-box
signals and small models; a dedicated reasoning verifier does better than both the posterior and the
cosine signal, consistent with our Lesson 2. None of these caveats touch the central, repeatedly
replicated finding: the connectivity boundary and the structure-not-Bayes equivalence.

\section{Conclusion}

Can Bayes help retrieval for LLM systems? Barely and sometimes---but the \emph{when} is sharp and
useful. It helps through structure (a graph prior, when evidence is genuinely connected and the
corpus is hard for similarity) and through adaptation (online updating, when the workload repeats). It
does not help as an uncertainty or decision signal, where one-line baselines win. And where it does
help, a cheap heuristic usually suffices, with the full posterior reserved for the one job only it does
well: gracefully gating structure across heterogeneous queries. For a field reaching reflexively for
Bayesian and uncertainty language around retrieval, the disciplined answer is to inject the right
structure, adapt when you can, and calibrate decisions with simple, transparent signals.

\section*{Reproducibility}
All experiments are scripted and use cached data; the scripts (\texttt{s1a\_graphgp.py},
\texttt{s2\_beaver.py}, \texttt{s3\_sql\_singlehop.py}, \texttt{s3\_sql\_hopgate.py},
\texttt{s3\_scifact.py}, \texttt{s4\_hotpot.py}, \texttt{s4\_diversity\_uq.py},
\texttt{s4\_adaptive.py}, \texttt{s5\_twowiki.py}, \texttt{s\_adapt\_beaver.py},
\texttt{downstream\_ex.py}) regenerate every number. API use was safe-by-default (dry-run with cost
estimate; explicit \texttt{--run}; call caps) and totaled $\approx\$6$.

\bibliographystyle{plainnat}
\bibliography{tas_refs}

\end{document}
